% !PRINTED 456
\weberpage{456}

Sind $\beta, \gamma \ldots$ Wurzeln von anderen Gleichungen in $\Omega$ oder auch andere Wurzeln derselben Gleichung, so erhalten wir durch gleichzeitige Adjunction von $\alpha, \beta, \gamma \ldots$ gleichfalls algebraische Körper über $\Omega$, die wir mit
\[ \Omega(\alpha, \beta, \gamma \ldots) \]
bezeichnen. Wir werden gleich sehen, dass diese Erweiterung des Begriffes algebraischer Körper nur eine scheinbare ist, und bleiben also zunächst bei der Adjunction \gesperrt{einer} Grösse $\alpha$, also bei der Betrachtung von $\Omega(\alpha)$ stehen.

Die Gleichung (1), deren Wurzel $\alpha$ ist, kann reducibel sein. Unter den irreducibeln Factoren von $F(x)$ ist aber wenigstens einer, der für $x = \alpha$ verschwindet, den wir mit $f(x)$ bezeichnen wollen; diese Function $f(x)$ ist, wenn wir den Co\"efficienten der höchsten Potenz von $x$ gleich 1 annehmen, völlig bestimmt, weil nach dem vorigen Paragraphen $x = \alpha$ nicht die Wurzel von zwei verschiedenen irreducibeln Gleichungen sein kann.

Die Gleichung $f(x) = 0$ hat also die Form
\weq{2}{x^n + a_1x^{n-1} + a_2x^{n-2} + \cdots + a_{n-1}x + a_n = 0,}
worin $a_1, a_2 \ldots a_n$ Grössen in $\Omega$ sind. Ist $n$ der Grad dieser Gleichung, so nennen wir $n$ auch den \gesperrt{Grad des Körpers} $\Omega(\alpha)$.

Alle Grössen $\Theta$ des Körpers $\Omega(\alpha)$ lassen sich ableiten durch Addition, Subtraction, Multiplication und Division aus $\alpha$ und aus Zahlen in $\Omega$; sie lassen sich also darstellen als rationale Functionen von $\alpha$ mit Co\"efficienten in $\Omega$, oder in der Form
\weq{3}{\Theta = \dfrac{\chi(\alpha)}{\psi(\alpha)},}
worin $\chi(x)$ und $\psi(x)$ Functionen in $\Omega$ sind. Da aber $\psi(\alpha)$ von Null verschieden sein muss, so ist $\psi(x)$ nicht durch $f(x)$ theilbar, und da $f(x)$ irreducibel angenommen ist, so ist $\psi(x)$ relativ prim zu $f(x)$.

Danach können wir nach \S.\,6, Satz II die Functionen $Q(x)$, $\varphi(x)$ in $\Omega$, und zwar $\varphi(x)$ höchstens vom Grade $n - 1$, so bestimmen, dass
\weq{4}{Q(x) f(x) + \varphi(x) \psi(x) = \chi(x)}
wird, und wenn wir also hierin $x = \alpha$ setzen, so dass $f(\alpha) = 0$ wird, so folgt
\weq{5}{\dfrac{\chi(\alpha)}{\psi(\alpha)} = \varphi(\alpha).}
